% Options for packages loaded elsewhere
\PassOptionsToPackage{unicode}{hyperref}
\PassOptionsToPackage{hyphens}{url}
\PassOptionsToPackage{dvipsnames,svgnames,x11names}{xcolor}
%
\documentclass[
  10pt,
]{article}
\usepackage{amsmath,amssymb}
\usepackage{iftex}
\ifPDFTeX
  \usepackage[T1]{fontenc}
  \usepackage[utf8]{inputenc}
  \usepackage{textcomp} % provide euro and other symbols
\else % if luatex or xetex
  \usepackage{unicode-math} % this also loads fontspec
  \defaultfontfeatures{Scale=MatchLowercase}
  \defaultfontfeatures[\rmfamily]{Ligatures=TeX,Scale=1}
\fi
\usepackage{lmodern}
\ifPDFTeX\else
  % xetex/luatex font selection
  \setmainfont[]{DejaVu Serif}
  \setsansfont[]{DejaVu Sans}
  \setmonofont[]{DejaVu Sans Mono}
  \setmathfont[]{Latin Modern Math}
\fi
% Use upquote if available, for straight quotes in verbatim environments
\IfFileExists{upquote.sty}{\usepackage{upquote}}{}
\IfFileExists{microtype.sty}{% use microtype if available
  \usepackage[]{microtype}
  \UseMicrotypeSet[protrusion]{basicmath} % disable protrusion for tt fonts
}{}
\makeatletter
\@ifundefined{KOMAClassName}{% if non-KOMA class
  \IfFileExists{parskip.sty}{%
    \usepackage{parskip}
  }{% else
    \setlength{\parindent}{0pt}
    \setlength{\parskip}{6pt plus 2pt minus 1pt}}
}{% if KOMA class
  \KOMAoptions{parskip=half}}
\makeatother
\usepackage{xcolor}
\usepackage[margin=2.0cm]{geometry}
\usepackage{longtable,booktabs,array}
\usepackage{calc} % for calculating minipage widths
% Correct order of tables after \paragraph or \subparagraph
\usepackage{etoolbox}
\makeatletter
\patchcmd\longtable{\par}{\if@noskipsec\mbox{}\fi\par}{}{}
\makeatother
% Allow footnotes in longtable head/foot
\IfFileExists{footnotehyper.sty}{\usepackage{footnotehyper}}{\usepackage{footnote}}
\makesavenoteenv{longtable}
\setlength{\emergencystretch}{3em} % prevent overfull lines
\providecommand{\tightlist}{%
  \setlength{\itemsep}{0pt}\setlength{\parskip}{0pt}}
\setcounter{secnumdepth}{-\maxdimen} % remove section numbering
% doc_header.tex -- LaTeX preamble additions used by build_docs.sh for every
% Markdown -> LaTeX -> PDF conversion (supplements, README, submission files).
\usepackage{etoolbox}
% longtable may not be loaded by pandoc when a document has no table
\usepackage{longtable}
% bold vectors: \bm -> \symbf under unicode-math (XeLaTeX), bm package otherwise
\makeatletter
\@ifpackageloaded{unicode-math}{\providecommand{\bm}[1]{\symbf{#1}}}{\usepackage{bm}}
\makeatother
% smaller type in long tables so that wide audit tables fit the text width;
% inside a table a line break is allowed after an underscore of plain text
% (pandoc writes it as \_), so that long identifiers wrap inside their column
% rows are set a little apart, so that symbols with indices do not touch
\AtBeginEnvironment{longtable}{\footnotesize\renewcommand{\arraystretch}{1.45}%
  \let\sptunderscore\_\def\_{\sptunderscore\allowbreak}}
\setlength{\LTcapwidth}{\textwidth}
% allow line breaks after slashes and hyphens inside table cells
\usepackage{xurl}
% code blocks: break over-long lines instead of running into the margin.
% Plain ``` blocks are typeset by pandoc as {verbatim}, highlighted ones
% (```bash) as {Highlighting}; both are re-declared as breakable fancyvrb
% environments.  The continuation line is marked with a small hooked arrow.
\usepackage{fvextra}
\RecustomVerbatimEnvironment{verbatim}{Verbatim}{breaklines,breakanywhere,fontsize=\small}
\makeatletter
\@ifundefined{Highlighting}{}{%
  \RecustomVerbatimEnvironment{Highlighting}{Verbatim}{breaklines,breakanywhere,fontsize=\small,commandchars=\\\{\}}}
\makeatother
% inline code (pandoc writes it as \texttt{...}): never hyphenate -- an inserted
% hyphen could be read as part of a file name -- but allow a line break after
% an underscore, slash, comma, colon or equals sign, so that long paths and
% column lists wrap inside the text width and inside table columns.
% Inside a formula (e.g. the name of a notebook symbol in an equation) the same
% command must first leave math mode, which \text of amsmath does.
\ExplSyntaxOn
\tl_new:N \l__spt_code_tl
\cs_new_protected:Npn \__spt_texttt:n #1
  {
    \tl_set:Nn \l__spt_code_tl {#1}
    \regex_replace_all:nnN { ([\/\,\:\=]) } { \1 \c{allowbreak} } \l__spt_code_tl
    \regex_replace_all:nnN { \c{_} } { \c{_} \c{allowbreak} } \l__spt_code_tl
    \group_begin: \ttfamily \hyphenchar \font = -1 \scan_stop: \l__spt_code_tl \group_end:
  }
\RenewDocumentCommand \texttt { m }
  {
    \mode_if_math:TF { \text { \__spt_texttt:n {#1} } } { \__spt_texttt:n {#1} }
  }
\ExplSyntaxOff
% A figure is written in the Markdown source as an image paragraph followed by
% a caption paragraph ("**Figure S1.** ..."), so that the caption is also
% visible on GitHub.  An image that starts a paragraph is centered.
\usepackage{graphicx}% pandoc loads it only for documents that contain an image
\NewCommandCopy\sptincludegraphics\includegraphics
\RenewDocumentCommand\includegraphics{O{} m}{%
  \ifvmode\begin{center}\sptincludegraphics[#1]{#2}\end{center}%
  \else\sptincludegraphics[#1]{#2}\fi}
% keep figures where they are referenced
\usepackage{float}
\floatplacement{figure}{H}
% running header with the document title
\usepackage{fancyhdr}
\pagestyle{fancy}
\fancyhf{}
\fancyhead[L]{\small PROJECT spacetimes --- documentation and provenance}
\fancyhead[R]{\small\thepage}
\renewcommand{\headrulewidth}{0.4pt}
% 16x16 matrices: tighter columns so that a full matrix fits the text width
\setlength{\arraycolsep}{3.2pt}
% The documents of this project contain long tagged symbols, file paths and
% identifiers that cannot be hyphenated.  A paragraph is therefore allowed to
% have looser lines before TeX gives up and lets a line stick out of the margin.
\AtBeginDocument{\tolerance=2000 \emergencystretch=5em \hbadness=2000}
\ifLuaTeX
  \usepackage{selnolig}  % disable illegal ligatures
\fi
\IfFileExists{bookmark.sty}{\usepackage{bookmark}}{\usepackage{hyperref}}
\IfFileExists{xurl.sty}{\usepackage{xurl}}{} % add URL line breaks if available
\urlstyle{same}
\hypersetup{
  pdftitle={P11 --- Stages of PROJECT spacetimes that were paused on November 12, 2026},
  pdfauthor={Prepared for Patrick Lee Nash, Ph.D.},
  colorlinks=true,
  linkcolor={blue},
  filecolor={Maroon},
  citecolor={blue},
  urlcolor={blue},
  pdfcreator={LaTeX via pandoc}}

\title{P11 --- Stages of PROJECT spacetimes that were paused on November
12, 2026}
\usepackage{etoolbox}
\makeatletter
\providecommand{\subtitle}[1]{% add subtitle to \maketitle
  \apptocmd{\@title}{\par {\large #1 \par}}{}{}
}
\makeatother
\subtitle{Provenance required by Hard Rule 7: what is not finished, why,
and how to finish it}
\author{Prepared for Patrick Lee Nash, Ph.D.}
\date{November 12, 2026}

\begin{document}
\maketitle

\hypertarget{what-this-document-is-for}{%
\section{1. What this document is for}\label{what-this-document-is-for}}

On Thursday, November 12, 2026, at 15:20 Pacific Standard Time, the
author instructed that all running stages be paused and completed later,
and that a new piece of work be started in the folder
\texttt{Revision\_Deflation}. This document records, for the stages that
were interrupted, (i) what was running, (ii) what state it was in, (iii)
what is known already, and (iv) what has to be done to complete it.
Nothing in the stages listed here is to be taken as finished.

The state of the repository at the moment of the pause is the commit
\texttt{3368a72}. All documents, programs and data of that commit are
complete in themselves; the stages below were to \emph{check and refine}
them.

\hypertarget{paused-stage-a-the-independent-verification}{%
\section{2. Paused stage A: the independent
verification}\label{paused-stage-a-the-independent-verification}}

\textbf{The instruction.} The task specification asks: ``Check and
verify all of your previous results/conclusions {[}\ldots{]} and refine
and refactor your solution.''

\textbf{What was running.} Four independent verification passes had been
started after the commit \texttt{3368a72}. Each pass was given to a
separate worker that had not seen the work being produced, that could
read the repository but not change it, and that was to re-derive or
re-compute what it checked by its own means. Each was to deliver a
report with three parts: findings (error, gap or minor point, each with
its exact place, the evidence and a proposed correction), what was
verified and how, and what was not checked.

\begin{longtable}[]{@{}
  >{\raggedright\arraybackslash}p{(\columnwidth - 2\tabcolsep) * \real{0.0429}}
  >{\raggedright\arraybackslash}p{(\columnwidth - 2\tabcolsep) * \real{0.9571}}@{}}
\toprule\noalign{}
\begin{minipage}[b]{\linewidth}\raggedright
Pass
\end{minipage} & \begin{minipage}[b]{\linewidth}\raggedright
Scope
\end{minipage} \\
\midrule\noalign{}
\endhead
\bottomrule\noalign{}
\endlastfoot
A & The two sets of Dirac matrices, the groups they generate, the chiral
projectors, the Dirac equations, the classification of the fermion pairs
(``dirac matrices''; manuscripts, Secs. II to IV and Appendix A). \\
B & The Vorton-Scalar model, the exchange amplitudes, the gap equation,
the bound pair, the effective action of Cornwall, Jackiw and Tomboulis,
the Higgs discussion (Supplements S1, S2, S4; manuscripts, Secs. V to
X). \\
C & Cosmology and dark matter (Supplement S3; manuscripts, Secs. XI and
XII and Appendix D). \\
D & Audit of the documents: every number against the data files,
cross-references, reproduction of the data from a clean copy, the symbol
rules of P1, the layout, and the rule that the manuscripts do not
mention the tools with which they were prepared. \\
\end{longtable}

\textbf{State.} The passes had run for about eleven minutes when the
service that runs them was restarted. They were started again and, a few
minutes later, stopped at the author's instruction. \textbf{No pass
delivered a report. No finding was received.} Therefore nothing in the
project has been confirmed, and nothing has been refuted, by this stage.
The working files of the passes are in a temporary workspace, not in
this repository, and should be considered lost.

\textbf{To complete.} Run the four passes again with the scopes of the
table; evaluate every finding; correct every confirmed error in the
programs and in the documents; record each correction in
\texttt{provenance/P9\_error\_log}; write the result as
\texttt{provenance/P10\_independent\_verification} (\texttt{.md},
\texttt{.tex}, \texttt{.pdf}), with the findings, their resolution, and
what was not checked; add P10 to the list of documents in
\texttt{README.md}.

\hypertarget{paused-stage-b-fermions-bound-to-a-vortex-line-a-known-imprecision}{%
\section{3. Paused stage B: fermions bound to a vortex line (a known
imprecision)}\label{paused-stage-b-fermions-bound-to-a-vortex-line-a-known-imprecision}}

\textbf{The statement that is to be corrected.} Section V H of both
manuscripts and Sec. 8 of Supplement S1 say that the Dirac equation in
the background of a vortex line of the Vorton-Scalar with winding number
\(n\) has ``\(|n|\) normalizable solutions of zero energy that move
along the line at the speed of light'', and attribute this to Jackiw and
Rossi. The count \(|n|\) is correct for \emph{one four-component} Dirac
field. It is not the count for the sixteen-component fields of this
project. The statement is recorded as error E43 in
\texttt{provenance/P9\_error\_log}, with the status ``identified, not
yet corrected''.

\textbf{What the analysis gives so far.} In the sector of fields that do
not depend on the extra coordinates, the one-particle Hamiltonian in a
Vorton-Scalar background \(\Phi_{VS}=(\phi_1+i\phi_2)/\sqrt2\) is

\[\mathsf h=\alpha^jp_j+g_{VS}\,\phi_1\,\beta+g_{VS}\,\phi_2\,(i\beta\chi),\qquad j=1,2,3,\]

with the five anticommuting Hermitian matrices
\(\alpha^1,\alpha^2,\alpha^3,\beta, i\beta\chi\) of either spacetime.
For a straight vortex line along \(x3\) with one unit of winding,
\(\phi_1+i\phi_2\propto e^{i\varphi}\) in the plane \((x1,x2)\), the
solutions of zero energy that do not depend on \(x3\) are a common
radial function, which decays away from the line, times a constant
spinor \(u\) that satisfies the two conditions

\[(i\beta\alpha^1)\,u=-u,\qquad(\alpha^2\beta\chi)\,u=-u .\]

The two matrices are Hermitian, have square one and commute. The
prototype \texttt{provenance/paused/vortex\_prototype\_algebra.py}
evaluates these conditions with the matrices of
\texttt{data/gamma1644.txt} and \texttt{data/gamma1671.txt}. Its
results, the same for both spacetimes, all exact:

\begin{longtable}[]{@{}
  >{\raggedright\arraybackslash}p{(\columnwidth - 2\tabcolsep) * \real{0.6420}}
  >{\raggedright\arraybackslash}p{(\columnwidth - 2\tabcolsep) * \real{0.3580}}@{}}
\toprule\noalign{}
\begin{minipage}[b]{\linewidth}\raggedright
Quantity
\end{minipage} & \begin{minipage}[b]{\linewidth}\raggedright
Result
\end{minipage} \\
\midrule\noalign{}
\endhead
\bottomrule\noalign{}
\endlastfoot
Dimension of the space of spinors \(u\) that satisfy both conditions &
4 \\
The matrix \(i\alpha^1\alpha^2\alpha^3\chi\) (the product of the five
matrices) & commutes with all five, is Hermitian and traceless and has
square one: it is \(\pm P\), the generator of \(U(1)_E\) \\
\(\alpha^3\) on the space of the four solutions & equal to
\(i\alpha^1\alpha^2\alpha^3\chi\); eigenvalues \(+1,+1,-1,-1\) \\
The same for the opposite winding (second condition with \(+u\)) &
dimension 4; \(\alpha^3=-i\alpha^1\alpha^2\alpha^3\chi\) \\
\end{longtable}

Because the energy of such a solution with the wavenumber \(k_3\) along
the line is \(\alpha^3k_3\), two of the four solutions move in the
direction \(+x3\) and two in the direction \(-x3\), at the speed of
light; the two directions are distinguished by the sign of \(P\). For a
general winding number the count per four-component field is \(|n|\)
(Jackiw and Rossi; an index theorem), so that the sixteen-component
field has \(4|n|\) zero modes, \(2|n|\) in each direction.

A second prototype,
\texttt{provenance/paused/vortex\_prototype\_profile.py}, solves the
equation of the profile of the vortex,
\(f''+f'/\rho-n^2f/\rho^2=\tfrac12f(f^2-1)\) with
\(\rho=m_s\times\)(distance from the line), by two methods. For \(n=1\)
it gives the slope \(f'(0)=0.4123772\), which is the known value
\(0.58319/\sqrt2\) of the vortex of the Ginzburg-Pitaevskii equation,
and the constant \(0.0343\) in the energy per unit length
\(\pi v_{VS}^2[\ln(m_sR)+0.0343]\) inside a radius \(R\), which
corresponds to the known form \(\ln(1.464\,R/\xi)\) of that equation.

\textbf{What is not done.} These are prototype results. They are not
part of the Rust program, they have no unit tests, and the documents
have not been changed.

\textbf{To complete.} (1) Add a module \texttt{vortex.rs} and a command
\texttt{spacetimes\ vortex} to the Rust program: the algebra above for
both spacetimes; the profile for \(n=1,2\); the radial functions of the
zero modes; the count of normalizable solutions in each angular-momentum
channel; unit tests; data files; one figure. (2) Correct the two
passages, and say in Sec. XIII of the manuscripts that the result is now
computed. (3) Enter the new symbols in the registry of P1, extend P8,
complete entry E43 of P9. (4) Run the fixed point of the symbol checks
and rebuild all documents that changed.

\hypertarget{other-things-that-were-still-to-be-done}{%
\section{4. Other things that were still to be
done}\label{other-things-that-were-still-to-be-done}}

\begin{itemize}
\tightlist
\item
  The final report to the author, with three questions: whether the
  affiliation printed in the manuscripts (``Independent researcher'',
  taken from the manuscript of the project NCSS) is the one wanted;
  whether the reference code CLASS may be built and run on this computer
  (it was not; the comparison uses the spectra stored in the project
  NCSS); and a reminder that the journal asks authors to state how a
  manuscript was prepared, while the manuscripts, as instructed, contain
  no such statement.
\item
  The code was executed in the cloud workspace, not under macOS: the
  sandbox on the author's Mac has no Rust compiler.
  \texttt{provenance/P8\_code} gives the commands for macOS; they have
  not been run there.
\end{itemize}

\hypertarget{what-is-not-affected-by-the-pause}{%
\section{\texorpdfstring{5. What is \emph{not} affected by the
pause}{5. What is not affected by the pause}}\label{what-is-not-affected-by-the-pause}}

The work in \texttt{Revision\_Deflation} is a separate piece of work
with hypotheses of its own. It uses the programs and matrices of this
project, and its documents state their own results. It changes one file
outside its folder: the script \texttt{code/scripts/build\_docs.sh},
which until then built every Markdown file of the repository, now leaves
the folder \texttt{Revision\_Deflation} to that folder's own copy of the
script (five lines changed on November 12, 2026). Without this change,
building all documents of the main project would also have rebuilt the
documents of the revision with the wrong settings and would have
overwritten its two manuscripts. No document, program or data file of
the main project was changed. The pending correction E43 concerns vortex
lines, which the revision treats independently.

\end{document}
